\documentclass[conference]{IEEEtran}
\IEEEoverridecommandlockouts
% The preceding line is only needed to identify funding in the first footnote. If that is unneeded, please comment it out.

\usepackage{cite}
\usepackage{amsmath,amssymb,amsfonts}
\usepackage{graphicx}
\graphicspath{{../figures/}}
\usepackage{algorithm}
\usepackage{algpseudocode}
\usepackage{textcomp}
\usepackage{xcolor}
\usepackage{placeins}
\usepackage{booktabs}
\usepackage{url}
\usepackage[hidelinks]{hyperref}

\def\BibTeX{{\rm B\kern-.05em{\sc i\kern-.025em b}\kern-.08em
    T\kern-.1667em\lower.7ex\hbox{E}\kern-.125emX}}

\begin{document}

\title{Design of a Multimodal Pipeline for Accessibility in E-Learning Content\\
}

\author{\IEEEauthorblockN{Asqar Nurym}
\IEEEauthorblockA{\textit{School of Software Engineering} \\
\textit{Astana IT University}\\
Astana, Kazakhstan \\
nurym.asqar@gmail.com}
}

\maketitle

\begin{abstract}
Educational video is now a core medium in online learning, but accessibility support is often fragmented across isolated features such as static captions, transcripts, or external assistive tools. This paper presents a unified multimodal pipeline for accessibility in e-learning video. The system combines automatic speech recognition (ASR), visual scene indexing, and on-demand audio description in a single processing flow that emits reusable artifacts for an accessible web player. The player integrates keyboard-first controls, ARIA semantics, and word-by-word subtitle rendering for attention support. A bilingual corpus evaluation was conducted on 24 educational videos (12 English, 12 Russian) using two configurations: ASR-only baseline (B0) and full pipeline (B1). Across the corpus, ASR confidence remained stable (mean 0.942), while the full pipeline achieved a mean real-time factor (RTF) of 0.433 and median RTF of 0.424, with faster-than-realtime processing in all 24 videos and mean 15-second scene coverage of 91.7\%. Results highlight practical performance under resource constraints and demonstrate the robustness of the adaptive scene density approach across diverse e-learning content types.

\end{abstract}

\begin{IEEEkeywords}
Accessibility, e-learning, educational video, multimodal pipeline, automatic speech recognition, captions, audio description, cognitive load.
\end{IEEEkeywords}

\section{Introduction}
\IEEEPARstart{V}{ideo-centric} learning has become dominant in higher education and professional training, but recorded lectures and screen-captured lessons remain unevenly accessible across learner groups. Learners who are d/Deaf or hard of hearing require synchronized captions as a primary accommodation, learners with low vision require access to salient visual information through audio description, and many learners in online contexts benefit from additional cognitive supports such as pacing control and navigational structure.

Existing practice often treats these needs as separate product features. In contrast, this work frames accessibility as an end-to-end pipeline problem: one processing flow should produce synchronized artifacts that can be reused by the playback layer and downstream tools. This framing is especially relevant under constrained deployment settings where compute, API quota, and implementation complexity are practical limitations.

The research question is: \emph{Can a resource-constrained multimodal pipeline produce synchronized accessibility artifacts with stable timing quality and practical processing latency across bilingual educational video content without requiring manual scene-budgeting?}

The contributions of this paper are as follows:
\begin{enumerate}
\item A unified multimodal pipeline that jointly produces word-level subtitle artifacts, scene-indexed visual descriptions, and structured summary/chapter outputs.
\item An adaptive scene indexing strategy that ensures visual coverage in low-variance content (e.g., screencasts) without static budget constraints.
\item An on-demand audio-description interaction model that avoids modifying source video and uses caching for repeated low-latency playback.
\item A technical scenario-based benchmark on a balanced 24-video EN/RU corpus, substituting for human-subject studies by quantifying technical coverage and temporal precision.
\end{enumerate}

\section{Related Work}
\subsection{Captions and Learning Accessibility}
Captions are a broadly beneficial support beyond disability-only scenarios. Prior evidence shows improvements in attention, comprehension, and recall for diverse user groups \cite{Gernsbacher2015}. In educational contexts, automated subtitles can be viable under realistic conditions, but quality depends on timing and segmentation \cite{Malakul2023}. Work on partial and synchronized captions further emphasizes that transcript quality alone is insufficient without temporal alignment \cite{Mirzaei2018}. Accessibility requirements are also reflected in WCAG guidance \cite{WCAG21}. \textbf{Gap:} many implementations still expose captions as isolated layers rather than integrated artifacts tied to navigation and multimodal support.

\subsection{Cognitive Load and Neurodiversity in Video Learning}
Instructional-video effectiveness is shaped by pacing, signaling, and learner control \cite{Mayer2020,CastroAlonso2021}. Media-induced challenges in online video lectures, including pace and intelligibility, directly affect learning conditions \cite{Lange2020}. Cognitive-load measurement in educational video remains context-dependent \cite{Kruger2016}, and recent neurodiversity-focused evidence highlights barriers caused by transcription errors and limited control affordances \cite{LeCunff2024,Horlin2024}. \textbf{Gap:} literature supports control and structure, but practical systems rarely generate these supports (e.g., chapters and scene-linked summaries) automatically from raw video in a single pass.

\subsection{ASR and Timestamping Foundations}
Modern ASR has reached strong practical utility through deep-learning pipelines \cite{Malik2020} and large-scale weakly supervised models \cite{Radford2022}. Adjacent methods such as WhisperX and CrisperWhisper target timestamp precision through forced alignment and attention-based refinement \cite{Bain2023,Wagner2024}. \textbf{Gap:} timestamping advances are often presented as stand-alone transcription improvements rather than components inside a deployable accessibility pipeline that also handles visual scene understanding and interactive rendering.

\subsection{Visual Accessibility and Audio Description}
For learners with low vision, Audio Description (AD) provides a critical spoken narrative of visual events. Traditionally, AD is authored manually, making it expensive and rare in rapid-turnaround educational content. While automated video captioning has scaled, automated AD is often omitted or treated as a separate post-production pipeline. \textbf{Gap:} there is a lack of integrated pipelines that produce both precise textual captions and semantic visual descriptions simultaneously, offering AD on-demand without inflating the source media's runtime.

\section{System Design}
\subsection{Pipeline Architecture and Artifacts}
The proposed system follows a \emph{Hybrid Multimodal Architecture}, which selectively employs deep neural networks for cognitively intensive tasks while utilizing deterministic algorithms for high-throughput media processing. To optimize performance on consumer-grade hardware, the pipeline implements a \emph{Parallel Orchestration Layer} using a \texttt{ThreadPoolExecutor}. This allows the acoustic (ASR) and visual (Scene Detection) branches to execute concurrently, significantly reducing the overall Real-Time Factor (RTF). Given an input video, the system emits three artifact groups:
\begin{itemize}
\item \textbf{Subtitle artifacts} $A_{sub}$: word-level timestamps, segment timeline, and WebVTT export.
\item \textbf{Scene artifacts} $A_{scene}$: time-indexed visual scenes with text-aware natural-language descriptions.
\item \textbf{Structure artifacts} $A_{sum}$: summary points and chapter markers derived from scene descriptions.
\end{itemize}

\begin{figure*}[!t]
\centering
\includegraphics[width=0.95\textwidth]{fig_pipeline_arch.drawio.png}
\caption{Unified accessibility pipeline architecture. Offline preprocessing generates reusable artifacts; during playback, audio description is requested on demand and cached for subsequent calls.}
\label{fig:architecture}
\end{figure*}

\subsection{Interaction Layer and Artifact Rendering}
To support dynamic karaoke-style highlighting, the pipeline uses a midpoint-based word grouper rather than boundary-based assignment. By assigning each word to a subtitle segment if its temporal midpoint falls within the segment's boundaries, the system achieves 100\% word coverage with zero temporal overlaps on the evaluated corpus.

Native HTML5 \texttt{<track>} elements lack the styling affordances required for granular, stateful word-by-word highlighting (e.g., current, previous, spoken, and upcoming states). Therefore, the web player bypasses native rendering, instead utilizing \texttt{requestAnimationFrame} and CSS custom properties for high-performance, 60fps subtitle overlays that automatically adapt to user preferences such as OpenDyslexic typography and High-Contrast modes.

Audio description is invoked on demand at playback time via Algorithm~\ref{alg:ad}: (i) the player resolves the nearest indexed scene within a fixed tolerance window (30 seconds), (ii) returns cached audio if present, otherwise (iii) synthesizes speech from the stored description and persists the result for future reuse, and (iv) pauses and resumes video playback around AD audio.

\begin{algorithm}[!t]
\caption{On-Demand Audio Description Resolution}
\label{alg:ad}
\begin{algorithmic}[1]
\Require $T_{\text{req}}$ (current playback time in seconds)
\Require $I_{\text{scene}}$ (array of indexed scenes)
\State $S_{\text{nearest}} \gets \text{find\_closest}(I_{\text{scene}}, T_{\text{req}})$
\If{$|S_{\text{nearest}}.time - T_{\text{req}}| > 30.0$}
    \State \Return \text{Error (No scene nearby)}
\EndIf
\State $CacheKey \gets S_{\text{nearest}}.id$
\If{\text{exists}($CacheKey$)}
    \State $Audio \gets \text{load\_from\_disk}(CacheKey)$
\Else
    \If{\text{google\_tts\_available}()}
        \State $Audio \gets \text{GoogleCloudTTS}(S_{\text{nearest}}.desc)$
    \Else
        \State $Audio \gets \text{EdgeTTS}(S_{\text{nearest}}.desc)$
    \EndIf
    \State $\text{save\_to\_disk}(CacheKey, Audio)$
\EndIf
\State $\text{pause\_video}()$
\State $\text{play\_audio}(Audio)$
\State $\text{resume\_video}()$
\State \Return $Audio$
\end{algorithmic}
\end{algorithm}

\section{Methodology}
\subsection{Hybrid Multimodal Architecture}
The proposed system follows a \emph{Hybrid Multimodal Architecture}, which selectively employs specialized pretrained and hosted components for cognitively intensive tasks while utilizing deterministic algorithms for high-throughput media processing. This decoupling ensures system stability and reduces the risk of error propagation across the pipeline.

\subsection{Bilingual Multimodal Processing}
The pipeline is designed as a language-agnostic framework with specific adapters for English and Russian.
\begin{enumerate}
\item \textbf{Acoustic Modeling:} We employ a large-scale weakly supervised model (Whisper) optimized for the \texttt{int8\_float16} compute type \cite{Radford2022}, achieving a balance between transcription accuracy and Real-Time Factor (RTF).
\item \textbf{Visual Semantic Extraction:} Visual changes are detected using a content-aware thresholding algorithm implemented via the PySceneDetect library \cite{PySceneDetect}. For each identified scene, a hosted multimodal description service generates a semantic description that encapsulates both textual (OCR-free) and structural visual data.
\item \textbf{Adaptive Scene Indexing:} To handle low-variance content such as screencasts, the system employs an \emph{Adaptive Uniform Sampling} fallback, ensuring a minimum scene density of two descriptions per minute when heuristic detection triggers are insufficient.
\end{enumerate}

\subsection{On-Demand Interaction Model}
The interaction model is centered on \emph{Just-in-Time (JIT) Synthesis}. Instead of pre-rendering audio descriptions for the entire video—which increases storage overhead and complicates synchronization—descriptions are synthesized and cached only upon explicit user request. This pattern supports personalized accessibility without modifying the underlying bitstream. To facilitate integration and support an \emph{API-first} design, the system provides a comprehensive OpenAPI specification for all interaction endpoints.

\section{Technical Evaluation}
\subsection{Experimental Design and Corpus Benchmark}
We evaluated the pipeline using a \emph{Bilingual Multimodal Corpus Benchmark} consisting of 24 educational videos (12 EN, 12 RU). The corpus was stratified by duration (short, medium, long) and content type (talking-head, slide-centric, screencast, and practical demo). 

The current evaluation utilizes an \emph{Adaptive Scene Density} strategy without a static scene cap, allowing the pipeline to scale with the visual complexity of the content. Processing was conducted on a localized workstation (AMD Ryzen 9, NVIDIA RTX 3050) to establish a baseline for consumer-grade deployability.

\begin{table}[t]
\caption{Corpus Composition}
\label{tab:corpus_comp}
\centering
\footnotesize
\setlength{\tabcolsep}{5pt}
\begin{tabular}{lccc}
\toprule
\textbf{Language} & \textbf{Short} & \textbf{Medium} & \textbf{Long} \\
\midrule
English & 4 & 4 & 4 \\
Russian & 4 & 4 & 4 \\
\bottomrule
\end{tabular}
\end{table}

Two configurations were compared:
\begin{itemize}
\item \textbf{B0 (ASR-only):} transcript and subtitle artifacts without visual branch.
\item \textbf{B1 (Full):} ASR + visual scene indexing + multimodal descriptions + summary/chapter generation.
\end{itemize}

\begin{table}[t]
\caption{Key Experimental Parameters}
\label{tab:exp_params}
\centering
\footnotesize
\setlength{\tabcolsep}{4.5pt}
\begin{tabular}{ll}
\toprule
\textbf{Component} & \textbf{Setting} \\
\midrule
ASR model & faster-whisper medium \\
ASR compute type & int8\_float16 (GPU) \\
Scene detector & PySceneDetect ContentDetector \\
Scene threshold & 27 \\
Scene indexing & Adaptive density (no static cap) \\
Description stage & Hosted multimodal API \\
TTS & Google Cloud Neural TTS (edge-tts fallback) \\
Pipeline mode & On-demand audio description \\
\bottomrule
\end{tabular}
\end{table}

\subsection{Formal Metrics}
The study reports engineering proxy metrics:
\begin{equation}
\text{low\_conf\_ratio} = \frac{w_{<0.5}}{w_{\text{all}}}
\end{equation}
\begin{equation}
\text{coverage}_{15s} = \frac{s_{\text{covered}}}{s_{\text{video}}}
\end{equation}
\begin{equation}
\text{RTF} = \frac{T_{\text{processing}}}{T_{\text{video}}}
\end{equation}
where RTF $< 1$ indicates faster-than-realtime processing.

For implementation, both $T_{\text{video}}$ and $s_{\text{video}}$ are measured as the effective processed span from the generated timeline end time, which keeps B0/B1 denominators consistent for each video.

Additional diagnostics are used to characterize alignment stability and visual indexing behavior:
\begin{equation}
\text{overlap\_ratio} = \frac{n_{\text{overlap}}}{\max(w_{\text{all}}-1,1)}
\end{equation}
\begin{equation}
\text{scene\_density} = \frac{n_{\text{scene}}}{T_{\text{video}}/60}
\end{equation}
\begin{equation}
\text{tail\_uncovered} = \max(0, T_{\text{video}} - t_{\text{last\_scene}})
\end{equation}
where $n_{\text{overlap}}$ is the number of overlapping adjacent word timestamps, and $t_{\text{last\_scene}}$ is the timestamp of the final indexed scene.

\subsection{Automation and Reproducibility Protocol}
All runs are executed by a single evaluation driver that reads a corpus manifest and processes each video in both configurations (B0 and B1). The driver stores per-video outputs and aggregate summaries in machine-readable files, including baseline comparisons and a JSON report for manuscript extraction. To reduce stale results, the protocol includes consistency checks for missing rows, non-positive runtimes, and out-of-range coverage values before aggregate statistics are generated. When valid artifacts already exist for a video/configuration, the driver reuses cached results, matching practical batch-processing behavior.

This setup allows recomputation of all reported tables from the same corpus specification and parameter set, while preserving practical caching behavior in the processing pipeline.

\subsection{Qualitative Scenario-Based Evaluation}
To investigate the system's efficacy beyond numerical proxies, we conducted a \emph{Scenario-Based Qualitative Analysis} of two contrasting content types:
\begin{enumerate}
\item \textbf{Scenario A (Slide-Centric Lecture):} In a 15-minute EN mathematics lecture, the system identified 42 visual transitions. The multimodal description stage correctly extracted LaTeX-style formulas (e.g., "$E=mc^2$") without a dedicated OCR engine, and the summary stage successfully synthesized these into three high-level chapter markers (Introduction, Core Proof, Summary). This demonstrates the pipeline's ability to maintain semantic continuity in text-dense educational material.
\item \textbf{Scenario B (Low-Variance Screencast):} A 5-minute RU software tutorial exhibited minimal visual motion, initially triggering only two heuristic scene changes. The \emph{Adaptive Uniform Sampling} fallback successfully injected periodic descriptions every 30 seconds, providing the blind or low-vision user with periodic "Visual Anchor Points" that would otherwise be missing in a standard detection-only pipeline.
\end{enumerate}

The results of these scenarios suggest that the unified architecture effectively compensates for the limitations of individual media-processing stages, ensuring a consistent accessibility layer across diverse e-learning formats.

\FloatBarrier
\section{Results and Analysis}
\subsection{Aggregate Performance}
Table~\ref{tab:lang_overall} reports language-level and overall aggregates. The ASR subsystem remains stable across languages (overall confidence mean 0.942; low-confidence ratio mean 0.026). Baseline B0 is consistently fast (overall mean RTF 0.193), while B1 remains practical (overall mean 0.433, median 0.424).

\begin{table}[t]
\caption{Aggregate Results by Language and Overall}
\label{tab:lang_overall}
\centering
\footnotesize
\setlength{\tabcolsep}{4.2pt}
\begin{tabular}{lcccccc}
\toprule
\textbf{Scope} & \textbf{N} & \textbf{B0} & \textbf{B1} & \textbf{B1} & \textbf{ASR} & \textbf{Cov.} \\
 &  & \textbf{RTF}\textsubscript{mean} & \textbf{RTF}\textsubscript{mean} & \textbf{RTF}\textsubscript{med} & \textbf{conf.} & \textbf{15s} \\
\midrule
EN & 12 & 0.158 & 0.369 & 0.353 & 0.948 & 91.58\% \\
RU & 12 & 0.228 & 0.497 & 0.504 & 0.937 & 91.82\% \\
Overall & 24 & 0.193 & 0.433 & 0.424 & 0.942 & 91.70\% \\
\bottomrule
\end{tabular}
\end{table}

\subsection{Baseline Overhead and Throughput Behavior}
Table~\ref{tab:duration_baseline} shows B0/B1 behavior by duration bucket. The B1-to-B0 mean ratio is highest on short videos, where fixed initialization overhead is less amortized. Table~\ref{tab:stage_latency} further breaks down the absolute processing time (in seconds) into the audio-only pipeline (B0) and the visual multimodal overhead. The visual overhead (scene detection, hosted description-service calls, and summary generation) scales with both the scene count and video duration.

\begin{table}[t]
\caption{Stage-Level Latency Breakdown (Mean Seconds)}
\label{tab:stage_latency}
\centering
\footnotesize
\setlength{\tabcolsep}{5pt}
\begin{tabular}{lccc}
\toprule
\textbf{Duration} & \textbf{Audio Pipeline} & \textbf{Visual Overhead} & \textbf{Total (B1)} \\
\midrule
Short & 28.8 & 49.0 & 77.7 \\
Medium & 69.4 & 74.2 & 143.6 \\
Long & 191.6 & 193.6 & 385.3 \\
\bottomrule
\end{tabular}
\end{table}

\begin{table}[t]
\caption{Baseline Comparison by Duration Bucket}
\label{tab:duration_baseline}
\centering
\footnotesize
\setlength{\tabcolsep}{4.2pt}
\begin{tabular}{lcccc}
\toprule
\textbf{Duration} & \textbf{N} & \textbf{B0 RTF}\textsubscript{mean} & \textbf{B1 RTF}\textsubscript{mean} & \textbf{B1/B0} \\
\midrule
Short & 8 & 0.192 & 0.522 & 2.71$\times$ \\
Medium & 8 & 0.178 & 0.365 & 2.05$\times$ \\
Long & 8 & 0.208 & 0.414 & 1.99$\times$ \\
Overall & 24 & 0.193 & 0.433 & 2.25$\times$ \\
\bottomrule
\end{tabular}
\end{table}

Figure~\ref{fig:rtf} visualizes B0 and B1 RTF across content types using boxplots (median line, interquartile box, and whiskers). Full-pipeline processing remains below realtime for all 24 videos in the corpus.

\begin{figure}[!t]
\centering
\includegraphics[width=\columnwidth]{fig_rtf_comparison.png}
\caption{Processing Speed (Real-Time Factor) across content types. The box spans the interquartile range (IQR), the horizontal line marks the median, and the whiskers denote the rest of the distribution up to $1.5 \times \text{IQR}$. A lower RTF indicates faster processing (the red dashed line represents realtime RTF=1.0).}
\label{fig:rtf}
\end{figure}

\subsection{Failure Cases and Practical Constraints}
Coverage variation is driven by visual characteristics. Because coding tutorials involve minor pixel changes (typing text) rather than global visual shifts, the raw detector can under-trigger; in one medium-length screencast, the initial pass returned zero scenes (0\% coverage). To mitigate this, we implemented a fallback mechanism: if fewer than 3 scenes are detected, the system applies uniform temporal sampling (one frame every 30 seconds). After this fallback, the same screencast reached 99.1\% coverage. With adaptive scene density (no static cap), all 24 videos achieved non-zero coverage, and the overall mean coverage reached 91.7\%. Figure~\ref{fig:coverage} shows the resulting content-type coverage distribution, illustrating high mean coverage (often $>87\%$) with error bars indicating the 95\% confidence intervals across the videos.

\begin{table}[t]
\caption{Content-Type Sensitivity (B1)}
\label{tab:content_sensitivity}
\centering
\footnotesize
\setlength{\tabcolsep}{4.2pt}
\begin{tabular}{lccccc}
\toprule
\textbf{Content type} & \textbf{N} & \textbf{B1 RTF} & \textbf{Cov. 15s} & \textbf{Scenes} & \textbf{Zero} \\
 &  & \textbf{mean} & \textbf{mean} & \textbf{mean} & \textbf{scene} \\
\midrule
Talking-head & 6 & 0.414 & 90.37\% & 29.83 & 0 \\
Slide-centric & 6 & 0.425 & 87.28\% & 33.67 & 0 \\
Screencast & 6 & 0.445 & 94.58\% & 39.00 & 0 \\
Practical demo & 6 & 0.449 & 94.55\% & 51.83 & 0 \\
\bottomrule
\end{tabular}
\end{table}

Table~\ref{tab:content_sensitivity} shows that screencasts achieve the highest aggregate coverage (94.58\%) once fallback sampling is enabled. The lowest mean coverage appears in slide-centric material (87.28\%), where slide transitions can be infrequent in some lectures. Practical demonstrations are the densest in detected scenes (mean 51.83 scenes) and exhibit slightly higher average B1 runtime due to richer visual variation and longer generated descriptions. With adaptive density, all content types achieve coverage above 87\%.

\begin{figure}[!t]
\centering
\includegraphics[width=\columnwidth]{fig_coverage.png}
\caption{Mean coverage within 15 s of indexed scenes by content type (95\% confidence intervals shown as error bars).}
\label{fig:coverage}
\end{figure}

\subsection{Interactive Latency and TTS}
Deferring Text-to-Speech (TTS) generation to playback time significantly reduces offline preprocessing costs, but requires careful latency management. We evaluated on-demand TTS with a small five-sentence micro-benchmark using Google Cloud Neural TTS. Cold-start requests (first synthesis) yielded a median latency of 765.3~ms. Cached requests yielded a median local file-hit latency of 0.05~ms. The cached value reflects storage lookup latency rather than a full network round-trip, but still supports the practical benefit of reuse without re-encoding source media.

\FloatBarrier
\section{Discussion}
\subsection{Interpretation of Baseline Overhead}
The B0/B1 comparison should be interpreted as an ablation-style systems analysis rather than a competitive benchmark against external platforms. B0 captures the minimum viable subtitle pipeline under the same ASR stack and hardware, while B1 adds visual indexing and multimodal generation. The measured overhead therefore estimates the incremental cost of accessibility functions beyond captions. The overall B1/B0 ratio of 2.25$\times$ indicates that the full multimodal pipeline adds approximately one additional unit of processing time relative to ASR-only processing.

The duration-stratified results suggest a predictable amortization effect: fixed initialization costs and API-call setup dominate short videos, while longer content allows overhead to be distributed over larger playback duration. This explains why B1/B0 ratios decrease from short to long buckets even when absolute B1 processing time increases.

\subsection{Deployment Implications}
From a deployment perspective, the results support a tiered operation model. When resource pressure is high, B0 can be used as a fast fallback for caption-first accessibility. When quality of multimodal support is prioritized, B1 remains practical for offline preprocessing and near-realtime batch throughput on consumer hardware. For interactive playback, on-demand description with caching reduces repeated latency without requiring video re-rendering.

The main operational bottleneck remains the hosted multimodal service. This dependency is manageable in prototype settings. The current adaptive scene density approach eliminates the need for a static scene cap, as demonstrated by the corpus results where scene counts naturally scale with content complexity (from 6 scenes in a short slide-centric video to 123 scenes in a long practical demo) while maintaining faster-than-realtime processing across all 24 videos.

\section{Threats to Validity}
\textbf{External validity:} although balanced, the 24-video corpus does not represent all real-world variability (e.g., heavy overlap speech, severe noise, extreme accents, non-academic genres).

\textbf{Construct validity:} reported metrics are engineering proxies (timing quality, coverage, throughput), not direct measures of learning gains or cognitive outcomes from human studies.

\textbf{Internal validity:} the multimodal description stage uses a hosted API model in a resource-constrained setup. This introduces network, rate-limit, and model-drift dependencies. In addition, part of corpus-level timing results reflects cache-aware reuse in repeated batch runs. Mitigation includes artifact persistence, caching, explicit stage-aware reporting, and transparent protocol disclosure.

\section{Conclusion}
This paper presented a unified multimodal pipeline for accessibility in e-learning video and evaluated it under resource constraints with a reproducible EN/RU corpus protocol. Compared with an ASR-only baseline (mean RTF 0.193), the full pipeline achieves a mean RTF of 0.433 with all 24 videos processed faster than realtime, while adding scene-aware accessibility functions unavailable in caption-only flows. The adaptive scene density approach yields a mean 15-second coverage of 91.7\% across all content types and duration buckets. The results support practical deployability for research prototypes and identify concrete next steps: domain-specific scene selection strategies, personalized summary granularity, and user studies to connect technical proxies with learning outcomes.

\bibliographystyle{IEEEtran}
\bibliography{references}

\end{document}
